\subsection{Effect of training budget and model size on optimal \moe sparsity}
\label{app:budget}

Recall that in ~\Cref{sec:effect_of_training_budget}, we demonstrated how the relationship between optimal total parameters $N*$, optimal active parameters $N*_a$, and optimal pretraining loss $L$ predictably changes as a function of sparsity $S$ and training budget $C$. 
In this section, we use the fitted isoFLOP surfaces to analyze how the optimal \moe sparsity $S^*$ changes as a function of total parameters $N$ and training budget $C$, as shown in~\Cref{fig:opt_sparsity}. Our main findings are:
\begin{itemize}[left=0pt,itemsep=2pt]
    \item Across all training budgets (ranging from 3e19 to 1e21 FLOPs), increasing the total parameters $N$ leads to an increase in the optimal sparsity level $S^*$. 
    \item For a fixed model size (i.e., total parameters $N$), increasing the training budget $C$ generally reduces the optimal sparsity level $S^*$. 
    \item The relationship between model size $N$ and optimal $S*$ is not linear. For smaller models (up to about $500 \cdot 10^6$  parameters), the optimal sparsity remains at $0$ (i.e., dense) for most compute budgets. 
\end{itemize}



% \newpage